% !TeX program = lualatex
% プレアンブル
\documentclass[11pt,a4paper]{article}
\usepackage{luatexja-fontspec}
\usepackage{amsmath,amssymb}
\usepackage{ntheorem}
\usepackage{mathtools}
\usepackage{stmaryrd}
\usepackage{enumitem}
\usepackage[authoryear,round]{natbib}
\usepackage{hyperref}
\usepackage{graphicx}
\usepackage{tikz}
\usetikzlibrary{arrows.meta,positioning,shapes,calc}
\usepackage{xcolor}
\usepackage{booktabs}
\usepackage{multirow}
\setlength{\parindent}{0pt}
\setlength{\parskip}{1\baselineskip}
\usepackage{bm}



\title{視点数付き寓におけるエルブラン関係と継続残余の構造\\
---二値截断の制限と四句分別の代数的定式化---}
\author{千葉将臣}
\date{2026-07-06}

\theoremheaderfont{\normalfont\bfseries}
\theorembodyfont{\normalfont}
\newtheorem{definition}{定義}[section]
\newtheorem{theorem}[definition]{定理}
\newtheorem{lemma}[definition]{補題}
\newtheorem{remark}[definition]{備考}
\newenvironment{proof}{\par\noindent\textbf{証明. }}{\hfill$\square$\par}

\begin{document}

\maketitle

\begin{abstract}
本論文では、圏論の関係代数的抽象化である「寓（Allegory）」\citep{FreydScedrov1990}に「視点数（PoV: Point of View）」および「継続残余（Continuation Residue）$\rho$」の概念を導入し、「視点数付き寓（PoV-Allegory）」として再構造化する。従来の寓は、二値截断部分圏 $B=\{0,1\}$ 上の静的な関係構造に依拠しており、自己言及的な無限ループや決定不能性に直面した場合に、評価を未定義または停止状態として扱う制約を有していた。本論文では、射の交わり（$\cap$）を単なる論理積ではなく「視点数の加算・併置」として再定義し、モジュラー律を「継続残余の動的分配則」として解釈する。これにより、計算機上のエルブラン関係の動的検証において発生する不完全性を、無記の不動点および空（Śūnyatā）の代数構造として型レベルで分離する枠組みを提示する。
\end{abstract}

\section{導入：静的関係の限界と二値截断}
従来の数理論理学、およびそれを圏論的に拡張した「寓（Allegory）」の体系は、ある関係 $R(x,y)$ が「成立するか（1）」「成立しないか（0）」という二値截断部分圏 $B$ を基礎としている。しかし、この枠組みを計算機上でのエルブラン解釈の動的検証（ランタイムの実行）に適用すると、自己言及的な無限退行や非停止性（未定義）に直面した際、体系は「証明不能」または「エラー」として評価を停止し、それ以降の状態を記述できない。

仏教哲学の文脈では、この状態は対象を「ある」「ない」の二値へ截断する記の問題として解釈できる。先行研究において著者は、この截断を回避し、意味と指示の差異を「継続残余 $\rho > 0$」として陽に記述する「無記（むき、avyākata）の意味論」を定式化した。

本論文の目的は、この視点数付きエルブラン関係の動的構造を、寓の関係代数的公理系と接続することにある。我々は、関係を単なる集合のインクルージョンではなく、観測者の「視点数（PoV）」によって層化された構造として定式化する。

\section{寓（Allegory）の標準的定義と二値的制約}
まず、寓の標準的定義を復習し、そこに含まれる二値的制約を明示する。

\begin{definition}[通常の寓]
寓 $\mathcal{A}$ とは、以下の条件を満たす圏である：
\begin{enumerate}
    \item 任意の対象 $X, Y$ に対し、ホムセット $\mathcal{A}(X,Y)$ は下半束（Semi-lattice）をなし、その交わりを $R \cap S$ と書く。
    \item 任意の射 $R: X \to Y$ に対し、逆元（Converse） $R^\circ: Y \to X$ が定義され、対合性（$(R^\circ)^\circ = R$）および単調性を満たす。
    \item 射の合成と交わり、および逆元は以下の\textbf{モジュラー律（Modular Law）}を満たす：
    \begin{equation}
        RS \cap T \subseteq R(S \cap R^\circ T)
    \end{equation}
\end{enumerate}
\end{definition}

通常の数学において、この公理系は $\mathbf{Set}$ 上の二項関係の圏 $\mathbf{Rel}$ を適切に抽象化する。しかし、ここでの包含関係 $\subseteq$ や等号 $=$ は、情報の同一化（円柱空間 $E_=$）を前提としており、無限遠点 $E_\infty$ における空間の非自明な同値構造（メビウス空間 $E_{\cong}$）を事前に除外している。有限の語彙と計算ステップに制約される計算機にとって、この除外は全域的な普遍性の構成的証明を困難にする要因となる。

\section{視点数付き寓（PoV-Allegory）の定式化}
この制約を扱うため、我々はホムセットの構造に観測文脈（層）を導入する。

\begin{definition}[真理値空間と残余空間]
真理値空間を四句分別（Caturkoṭi）に基づき、以下のように分離して定義する：
\begin{itemize}
    \item $\mathcal{C} = \{C_1(\text{是}), C_2(\text{非}), C_3(\text{亦是亦非}), C_4(\text{非是非非})\}$
    \item 残余空間 $\mathcal{L} = \mathbb{R}_{\ge 0}$ （継続残余 $\rho$ の大きさを示す）
    \item 視点数空間 $\mathcal{V} = \mathbb{N} \cup \{\infty\}$ （PoVのインデックス）
\end{itemize}
\end{definition}

\begin{definition}[視点数付き寓の射]
対象 $X, Y$ 間の射 $R: X \to_{\langle n, \rho \rangle} Y$ とは、各ペア $(x,y) \in X \times Y$ に対し、評価推移列 $\gamma$ を介して真理値、視点数、および継続残余の三つ組を割り当てる多値的な関係である。
\end{definition}

ここで、寓における基本演算である「交わり（$\cap$）」の演算を、単一階層上の論理積ではなく、視点数を増加させる演算として再定義する。

\begin{definition}[交わり演算とPoVの増加]
2つの射 $R, S: X \to Y$ の交わり $R \cap S$ とは、それぞれの観測文脈の「併置」である。
あるペア $(x,y)$ において、
$R$ が $PoV = n$, $S$ が $PoV = m$ で観測されているとき、その交わりの視点数は次のように定義される：
\begin{equation}
    \text{PoV}(R \cap S) = n + m
\end{equation}
これは、「$A$ かつ $B$」という結合が、計算過程では二つの独立した証明文脈の併置として現れ、その結果として必要な視点数（PoV）が増加することを意味する。

また、このとき評価文脈間の干渉が、以下の通り継続残余の増大として現れる：
\begin{equation}
    \rho(R \cap S) = \rho(R) + \rho(S) + \Delta_{n,m}
\end{equation}
\end{definition}

\begin{definition}[逆元と否定方向の評価]
射 $R$ の逆元 $R^\circ$ とは、ドメインとコドメインの反転を意味すると同時に、四句分別における $C_2$（非：否定方向の単一視点操作）への対応を与える。したがって、
\begin{equation}
    \text{PoV}(R^\circ) = \text{PoV}(R)
\end{equation}
であり、関係の向きのみが反転する。

\end{definition}

\section{モジュラー律の動的解釈：継続残余の分配方程式}
視点数付き寓において、寓の最大の特徴であるモジュラー律：
\begin{equation}
    RS \cap T \subseteq R(S \cap R^\circ T)
\end{equation}
は、もはや静的な集合の包含関係ではない。これは、左辺で生じる大域的な証明障害を、右辺における逆向き評価文脈の導入によって局所的に分配する残余の不等式として解釈できる。

\begin{theorem}[残余モジュラー定理]
左辺 $RS \cap T$ において発生する大域的な $E_\infty$ 型の証明障害を $\rho_{\text{left}}$ とし、右辺 $R(S \cap R^\circ T)$ において実行時に局所化される残余の総和を $\rho_{\text{right}}$ とするとき、以下の不等式が維持される：
\begin{equation}
    \rho_{\text{left}} \ge \rho_{\text{right}}
\end{equation}
すなわち、右辺の手続きは、大域的なループを形成する前に局所的な反転視点 $R^\circ$ を $T$ に導入するため、実行時に残余を事前に局所化できる。
\end{theorem}

\begin{proof}
千葉の『継続残余としての証明障害』の補題4.2より、動的検証における評価文脈の導入は、実行時の継続スタックを短縮する。右辺における $R^\circ T$ の交わりは、ターゲット空間 $T$ に含まれる非自明な構造を反転射によって補正するため、外側から改めて $R$ を合成した際の未完結性の累積を厳密に抑制する。したがって、残余は右辺において抑制される。
\end{proof}

\section{無記の不動点と空（Śūnyatā）への収束}
合成の反復が無限に続き、PoV が $\infty$ に達したとき、体系内の恒等射としての解釈は維持できず、通常の寓では評価が停止する。しかし、視点数付き寓においては、この極限は無記の不動点として扱われる。

\begin{theorem}[空への最大固定点収束]
\citet{Kleene1952} の有限主義的な計算観の下で実装された視点数付き寓において、任意の自己言及的射の無限列 $R \subseteq R^2 \subseteq R^3 \subseteq \dots$ は、真理値空間において $C_4$（非是非非、未定義の集積）へ向かうが、残余空間 $\mathcal{L}$ が極限消去（$E_\infty$ の包摂）を行うことで、最大固定点（dfix）としての $C_5$（空：Śūnyatā）へ収束する。
\end{theorem}

\begin{remark}
このとき、既存の数学が「矛盾」あるいは「計算不可能性」として排除してきた成分は、無記（$C_0$）として残余空間 $\mathcal{L}$ へ分類されるため、計算過程は論理的整合性を保ったまま後続の処理を継続できる。
\end{remark}

\section{結語}
本論文が提示した「視点数付き寓（PoV-Allegory）」は、既存の寓に視点数と継続残余を追加する拡張である。この拡張は、二値的な同一化を前提とする静的関係構造を、計算過程に即した関係構造として再定式化するものである。

本体系は、寓を有限主義的な計算観の下で関係化し、視点数を明示的なパラメータとして導入することにより、大域的な $E_\infty$ 型の問題を制御可能な構造として記述する。この体系は、定理証明支援系やキュービカル型理論において、空に対応する極限状態を操作的に扱うための基礎的枠組みとなり得る。


\begin{thebibliography}{9}

\bibitem[Kleene(1952)]{Kleene1952}
S.~C.~Kleene.
\newblock Introduction to Metamathematics.
\newblock North-Holland Publishing Company, Amsterdam, 1952.

\bibitem[Freyd and Scedrov(1990)]{FreydScedrov1990}
P.~J.~Freyd and A.~Scedrov.
\newblock Categories, Allegories.
\newblock North-Holland, Volume 39, 1st edition, 1990.

\end{thebibliography}

\end{document}